\documentclass[a4paper,12pt]{article}
\usepackage[utf8]{inputenc}
\usepackage[margin=1in]{geometry}
\usepackage{xcolor}
\usepackage{hyperref}
\usepackage{graphicx}
\usepackage{tcolorbox}
\usepackage{booktabs}
\usepackage{listings}

% Code styling
\definecolor{codegreen}{rgb}{0,0.6,0}
\definecolor{codegray}{rgb}{0.5,0.5,0.5}
\definecolor{codepurple}{rgb}{0.58,0,0.82}
\lstset{
    commentstyle=\color{codegreen},
    keywordstyle=\color{magenta},
    numberstyle=\tiny\color{codegray},
    stringstyle=\color{codepurple},
    basicstyle=\ttfamily\footnotesize,
    breaklines=true,
    captionpos=b,
    keepspaces=true,
    numbers=left,
    showspaces=false,
    showstringspaces=false,
    frame=single
}

\begin{document}

% ==========================================
% DEDICATED COVER PAGE
% ==========================================
\begin{titlepage}
    \centering
    \vspace*{2cm}
    
    {\Huge \textbf{Software Testing Mini Project} \par}
    \vspace{1.5cm}
    {\Large \textbf{My Journey with Mutation Testing: \\ Hardening the LocalCUA Agent} \par}
    
    \vspace{3.5cm}
    
    {\Large \textbf{Submitted To:} \par}
    {\Large Sir Usman Gous \par}
    
    \vspace{2.5cm}
    
    {\Large \textbf{Report By:} \par}
    {\Large Muhammad Abu Huraira \par}
    {\Large Roll Number: 22F-3853 \par}
    
    \vfill
    
    {\Large \textbf{Date:} \par}
    {\Large April 30, 2026 \par}
    
    \vspace{2cm}
\end{titlepage}

% ==========================================
% TABLE OF CONTENTS
% ==========================================
\tableofcontents
\newpage

% ==========================================
% MAIN DOCUMENT BODY
% ==========================================
\section{Introduction: The LocalCUA Architecture}
To understand why testing is so critical for this project, you first have to understand the architecture we built. The \textbf{LocalCUA Agent} is a system that allows users to control their computer using natural language commands. To make this magic happen, the system is divided into three major components:

\begin{enumerate}
    \item \textbf{The Flutter Front-End:} Think of this as the face of the application. It provides the user interface where you type your commands. Behind the scenes, it acts as our eyes, continuously capturing screenshots of your current window to understand what is happening on screen.
    \item \textbf{The VLM (Vision-Language Model):} This is a fine-tuned AI model that acts as the creative brain. It takes the screenshot from Flutter and your text command, looks at the UI, and decides what actions need to be taken (e.g., clicking a specific coordinate or typing a word).
    \item \textbf{The Middle Layer (Agent \& Flask/FastAPI):} This is the connective tissue that holds everything together. The Flutter front-end talks to the Flask backend, which coordinates with the VLM. The most critical part of this middle layer is the \textbf{Action Parser}. It takes the VLM's abstract text output and translates it into physical mathematical coordinates that the computer can actually click.
\end{enumerate}

\section{System Workflow: Running Locally with Ollama}
One of the most powerful features of the LocalCUA Agent is that it runs entirely locally, ensuring complete privacy and zero cloud API costs. Here is exactly how the components talk to each other:
\begin{itemize}
    \item \textbf{Step 1 (The Trigger):} You open the Flutter app and type a natural language command (e.g., "Click the submit button"). The Flutter front-end instantly takes a screenshot of your active window and packages your text command.
    \item \textbf{Step 2 (The Flask Bridge):} Flutter sends this data package via a local HTTP request to the Flask/FastAPI middle layer running on your machine.
    \item \textbf{Step 3 (Ollama Execution):} The Flask server forwards the screenshot and command to \textbf{Ollama}. Ollama is a local AI engine hosting our fine-tuned Vision-Language Model. Ollama runs the heavy AI inference completely offline using your local hardware and returns an abstract text response detailing what action needs to be taken.
    \item \textbf{Step 4 (The Action Parser):} Finally, the Flask middle layer receives Ollama's text, runs it through the Action Parser to calculate exact mathematical pixel coordinates, and physically moves and clicks your mouse.
\end{itemize}

\section{Why Testing the Middle Layer Matters}
If the Flutter UI glitches, a button might look ugly. If the VLM hallucinates, it might output a confusing sentence. But if the \textbf{Middle Layer (Action Parser)} fails, the consequences are disastrous. 

Imagine you ask the agent to click "Save." The VLM correctly identifies the save button. But if my Action Parser has a slight mathematical bug in how it resizes images or rounds coordinates, the mouse might land a few pixels to the left and click "Delete" instead. Because the middle layer actually executes physical actions on your computer, it \textbf{must} be flawless.

This is why I chose \textbf{Mutation Testing} for the Flask backend. Standard unit testing only checks "does the code run without crashing?" Mutation testing acts like an active saboteur. It secretly changes my logic (injecting "fake bugs" or mutants) to see if my test suite is actually paying attention. If my tests pass while the code is secretly broken, my tests are weak. If my tests fail, my tests are doing their job.

\section{Task 1: The Setup Journey (Enter WSL)}
My first challenge was just getting the mutation testing tool, \verb|mutmut|, to run smoothly. Because I develop on a Windows machine, the heavy file I/O operations of running hundreds of tests simultaneously led to nasty file-locking issues. The mutation runner would randomly crash because Windows wouldn't let it edit files fast enough.

To solve this, I moved the entire testing pipeline into \textbf{WSL (Windows Subsystem for Linux)}. In this Ubuntu environment, I set up a pristine testing pipeline:
\begin{itemize}
    \item Installed \verb|mutmut| and \verb|pytest|.
    \item Configured a \verb|setup.cfg| file to tell mutmut exactly which file to attack (\verb|action_parser.py|) and where to find my tests.
    \item Wrote a bash script (\verb|wsl_mutation_runner.sh|) that violently purges all caches before each run, ensuring I never get false positives.
\end{itemize}

\section{Task 2: The Baseline Reality Check}
Before mutation testing, I ran standard code coverage. The results were fantastic—over 90\% coverage! I felt like my code was bulletproof. Then, I let the mutmut saboteur loose. Here is what actually happened:

\begin{table}[h!]
\centering
\caption{The Baseline Mutation Run}
\begin{tabular}{|l|c|}
\hline
\textbf{Metric} & \textbf{Value} \\ \hline
Total Mutants (Fake Bugs) Generated & 225 \\ \hline
Mutants Killed (My tests noticed the bug) & 177 \\ \hline
Mutants Survived (My tests were blind) & \textbf{48} \\ \hline
Initial Mutation Score & $\sim 78.6\%$ \\ \hline
\end{tabular}
\end{table}

Seeing 48 surviving mutants was a massive reality check. It meant there were 48 distinct ways my Action Parser could be logically broken without my tests ever noticing. I decided to hunt down 6 of the sneakiest surviving mutants and write super-targeted, humanized tests to kill them.

\section{Task 3: The Mutant Hunt}
Here is a detailed look at exactly what these mutants were, how they tricked my tests, and the clever traps I had to build to defeat them.

\subsection{Mutant 1: The Aspect Ratio Boundary (Mutant 20)}
\begin{itemize}
    \item \textbf{The Sabotage:} The code checks if an image is too wide using \verb|max_ratio = 10|. Mutmut secretly changed the \verb|>| sign to a \verb|>=| sign.
    \item \textbf{Why it survived:} My tests only checked normal images (ratio 2.0) and super stretched images (ratio 20.0). Since 2.0 passes both operators and 20.0 fails both, the mutant blended right in. I had never tested an image that was \textit{exactly} at the limit.
    \item \textbf{The Defeat:} I wrote a test that passes an image of exactly 100x1000 pixels (a ratio of exactly 10.0). The original code allows this (\verb|10 > 10| is False), but the mutant blocks it (\verb|10 >= 10| is True). The test caught the discrepancy and killed the mutant.
\end{itemize}

\subsection{Mutant 2 \& 3: The Silent Log Assassins (Mutants 18 \& 21)}
\begin{itemize}
    \item \textbf{The Sabotage:} Mutmut injected the string \verb|XX| straight into my backend \verb|logger.warning| messages.
    \item \textbf{Why it survived:} Like many developers, my tests only cared about what the function returned. I didn't write a single test to read the system logs. Mutmut was actively corrupting my debugging logs, and my tests were entirely blind to it.
    \item \textbf{The Defeat:} I utilized Pytest's \verb|caplog| fixture. I forced the system to generate warnings, and then strictly asserted that the log text was perfectly formatted without any \verb|XX|. The mutants were instantly caught.
\end{itemize}

\subsection{Mutant 4: The Sneaky Quote Stripper (Mutant 75)}
\begin{itemize}
    \item \textbf{The Sabotage:} The parser strips quotes from the VLM's text using \verb|value.strip("'\"")|. Mutmut altered this to \verb|value.strip("XX'\"XX")|.
    \item \textbf{Why it survived:} Python's \verb|strip()| removes any matching characters from the outside edges of a word. Because my test inputs were simple words like \verb|"hello"|, there were no \verb|X|'s on the edges. Both the original code and the mutant stripped the quotes and left \verb|hello|.
    \item \textbf{The Defeat:} I had to outsmart it. I wrote a test passing the word \verb|"helloX"|. The original code strips the quotes and returns \verb|helloX|. The mutated code strips the quotes \textit{and} the X, returning just \verb|hello|. My test immediately noticed the missing 'X' and killed the mutant.
\end{itemize}

\subsection{Mutant 5: The List Offset Glitch (Mutant 105)}
\begin{itemize}
    \item \textbf{The Sabotage:} When splitting the VLM text to find the last action, the code grabs the final item using the list index \verb|[-1]|. Mutmut changed this index to \verb|[+1]|.
    \item \textbf{Why it survived:} In my tests, I usually only sent a single action from the VLM. When you split a string with one action, you get a list of two items. In Python, index \verb|[-1]| and index \verb|[1]| on a two-item list point to the exact same element! 
    \item \textbf{The Defeat:} I wrote a test that intentionally feeds the parser \textbf{two} different actions (a "click" followed by a "type"). This created a list of three items. \verb|[-1]| correctly grabbed the last item ("type"), while the mutated \verb|[+1]| grabbed the middle item ("click"). The test failed, and the mutant died.
\end{itemize}

\subsection{Mutant 6: The Fallback Inversion (Mutant 117)}
\begin{itemize}
    \item \textbf{The Sabotage:} A safety check \verb|if smart_resize is None:| was inverted to \verb|is not None:|.
    \item \textbf{Why it survived:} When the image resizer fails, the backend falls back to raw coordinates and logs a warning. Since my tests only verified that the parser didn't crash, the missing backup warning went completely unnoticed.
    \item \textbf{The Defeat:} I wrote a test that intentionally forced the resizer to fail, and then explicitly asserted that the exact fallback warning was emitted. Because the inverted mutant suppressed the warning, the test caught the sabotage.
\end{itemize}

\section{Task 4: The Final Results}
After implementing these 6 targeted "Killer Tests," I ran the WSL mutation pipeline one last time to measure the improvement.

\begin{table}[h!]
\centering
\caption{The Final Mutation Run}
\begin{tabular}{|l|c|c|c|}
\hline
\textbf{Metric} & \textbf{Before} & \textbf{After} & \textbf{Change} \\ \hline
Total Mutants Generated & 225 & 225 & 0 \\ \hline
Mutants Killed & 177 & 183 & \textbf{+6} \\ \hline
Mutants Survived & 48 & 42 & \textbf{-6} \\ \hline
\end{tabular}
\end{table}

% ==========================================
% DEDICATED OBJECTIVES & CONCLUSION
% ==========================================
\newpage
\section{Project Objectives Achieved}
Through this mini-project, the following core objectives were successfully met:
\begin{itemize}
    \item \textbf{Environment Isolation:} Successfully isolated a complex, resource-heavy mutation testing pipeline from Windows to WSL, eliminating destructive file-locking errors.
    \item \textbf{Vulnerability Discovery:} Identified 48 logical blind spots in the Action Parser that standard line-coverage tests failed to detect.
    \item \textbf{Targeted Eradication:} Designed and implemented 6 highly specific, humanized unit tests utilizing Pytest fixtures (like \verb|caplog|) and boundary condition math to explicitly trap and kill surviving mutants.
    \item \textbf{Quantitative Improvement:} Raised the total number of killed mutants from 177 to 183, directly hardening the physical mouse-clicking logic of the AI agent.
\end{itemize}

\section{Conclusion \& What I Have Learned}
This assignment completely reshaped how I view software testing. Before this project, I believed that achieving 100\% test coverage meant my code was bulletproof. Mutation testing shattered that illusion—it proved that standard tests only verify that code \textit{executes}, but they do not prove that code executes \textit{correctly} under tricky edge conditions. 

By looking at my Middle Layer architecture through the eyes of a saboteur, I discovered that I was blindly accepting edge cases, ignoring system logs, and getting lucky with list indices. Writing these simple but highly detailed tests didn't just improve a metric on a dashboard; it genuinely made the LocalCUA Agent safer, smarter, and far more reliable when interpreting VLM commands. I have learned that true software reliability requires testing not just the happy paths, but actively anticipating and guarding against malicious logical shifts.

\end{document}
